% 填写题目、基本信息及中英文摘要；空基本信息保留横线，不列班级。
% 材料齐全、内容核实后改为 \ReportDraftfalse，移除模板草稿提示。
\ReportDrafttrue
% 有原始材料或教师要求时改为 \ReportAppendixtrue，启用正文末尾附录。
\ReportAppendixfalse
\newcommand{\ReportTitle}{实验题目}
\newcommand{\ReportTitleEN}{English title to be supplied}
\newcommand{\ReportAuthor}{\ReportBlank{2.2cm}}
\newcommand{\ReportAuthorEN}{\ReportBlank{2.2cm}}
\newcommand{\ReportStudentID}{\ReportBlank{2.5cm}}
\newcommand{\ReportTeacher}{\ReportBlank{2.2cm}}
\newcommand{\ReportExperimentDate}{\ReportBlank{2.5cm}}
\newcommand{\ReportAbstractCN}{\ReportTODO{用100--200字概括实验目的、方法、主要结果和结论；有真实数据时包含至少一个关键数值及单位。缺少数据时明确说明，不能虚构结果。}}
\newcommand{\ReportKeywordsCN}{\ReportTODO{3--5个关键词，以分号分隔。}}
\newcommand{\ReportAbstractEN}{To be supplied: translate the Chinese abstract faithfully, including the same physical quantities, units, numerical results and conclusions.}
\newcommand{\ReportKeywordsEN}{To be supplied: 3--5 terms corresponding to the Chinese keywords.}
